\documentclass[11pt]{article}

\usepackage[margin=1in]{geometry}
\usepackage{amsmath,amssymb,amsthm}
\usepackage{graphicx}
\usepackage[hidelinks]{hyperref}

\newtheorem{theorem}{Theorem}
\newtheorem{corollary}{Corollary}
\newtheorem{lemma}{Lemma}
\newtheorem{remark}{Remark}

\newcommand{\N}{\mathbb N}
\newcommand{\R}{\mathbb R}
\newcommand{\E}{\mathbb E}
\newcommand{\Prob}{\mathbb P}
\newcommand{\ellinf}{\ell_\infty}
\newcommand{\ellone}{\ell_1}
\newcommand{\co}{c_0}
\newcommand{\Gammafin}{[\N]^{<\omega}}
\newcommand{\norm}[1]{\left\lVert #1\right\rVert}

\begin{document}

\begin{center}
{\Large A second-predual counterexample for Gaussian random-field samples}

\medskip
{\small Run \texttt{fa\_banach\_001}; source arXiv:2203.04650}
\end{center}

\section*{Source question}

Korolev, Latz, and Schoenlieb ask in the Outlook of \emph{Gaussian random
fields on non-separable Banach spaces} whether the presence of a second
predual can strengthen their bounded-pairing and weak-star-completion
conclusions to actual membership of the sample in \(U\).

\begin{figure}[h]
\centering
\includegraphics[width=\textwidth]{figures/open_problem_crop.png}
\caption{The source Outlook question, local build PDF page 14.}
\end{figure}

In the notation of the source paper, a positive nuclear covariance is written
as
\[
  c=\sum_{i=1}^\infty \lambda_i \phi_i\otimes \phi_i,
  \qquad \lambda_i\ge 0,\quad \sum_i\lambda_i<\infty,
\]
with \(\norm{\phi_i}_U=1\), and the associated formal Gaussian sample is
\[
  \theta=\sum_{i=1}^\infty \lambda_i^{1/2}\xi_i\phi_i,
\]
where the \(\xi_i\)'s are independent standard Gaussian variables.  The source
paper proves that each fixed predual pairing is finite almost surely and that
\(\theta\) belongs to a weak-star completion.  The question is whether, in
second-predual cases, one can conclude \(\theta\in U\) almost surely.

\section*{Proof intuition}

The obstruction is that scalar testing by each fixed predual vector does not
control a supremum over too many coordinates.  In \(\ell_\infty\), make the
coordinates correspond to finite subsets \(F\subset\mathbb N\).  The \(i\)-th
covariance direction is the indicator of the family of finite sets containing
\(i\).  Then the coordinate indexed by \(F\) is the finite partial subseries
\(\sum_{i\in F}\xi_i/i\).  Taking \(F\) to be the positive indices among
\(\{1,\ldots,N\}\), these coordinates dominate a Bernoulli-thinned harmonic
series, and hence become unbounded almost surely.

\section*{Counterexample}

\begin{theorem}
Let \(U=\ell_\infty(\Gammafin)\).  Then \(U=\co(\Gammafin)^{**}\), so \(U\)
has the second predual \(\co(\Gammafin)\).  There are norm-one vectors
\((u_i)\subset U\) and positive numbers \((\lambda_i)\in\ell_1\) such that
the positive nuclear kernel
\[
  c=\sum_{i=1}^\infty \lambda_i u_i\otimes u_i
      \in U\widehat{\otimes}_\pi U
\]
has the source paper's scalar bounded-pairing property, but the corresponding
Gaussian sample
\[
  \theta=\sum_{i=1}^\infty \lambda_i^{1/2}\xi_i u_i
\]
does not belong to \(U\) almost surely.
\end{theorem}

\begin{proof}
Let \(\Gammafin\) be the countable set of all finite subsets of \(\mathbb N\).
For \(i\in\mathbb N\), define \(u_i\in \ell_\infty(\Gammafin)\) by
\[
  u_i(F)=
  \begin{cases}
  1, & i\in F,\\
  0, & i\notin F,
  \end{cases}
  \qquad F\in\Gammafin.
\]
Then \(\norm{u_i}_\infty=1\).  Put \(\lambda_i=i^{-2}\).  Since
\(\sum_i\lambda_i<\infty\), the series
\[
  c=\sum_{i=1}^\infty i^{-2}u_i\otimes u_i
\]
converges in the projective tensor product.  It is symmetric and positive in
the covariance sense, because for every \(\eta\in\ell_1(\Gammafin)\),
\[
  \langle \eta,C\eta\rangle
  =\sum_{i=1}^\infty i^{-2}\langle \eta,u_i\rangle^2\ge 0.
\]

Let
\[
  \theta(F)=\sum_{i\in F}\frac{\xi_i}{i}, \qquad F\in\Gammafin.
\]
Each coordinate is a finite sum.  Moreover, if \(\eta\in\ell_1(\Gammafin)\)
and
\[
  b_i=\langle \eta,u_i\rangle=\sum_{F\ni i}\eta_F,
\]
then \(|b_i|\le \norm{\eta}_1\), so
\[
  \sum_{i=1}^\infty \frac{b_i^2}{i^2}
  \le \norm{\eta}_1^2\sum_{i=1}^\infty \frac{1}{i^2}<\infty.
\]
Hence the scalar Gaussian series
\[
  \sum_{i=1}^\infty \frac{\xi_i b_i}{i}
\]
converges almost surely.  This is exactly the bounded-pairing conclusion for
each fixed predual vector \(\eta\).

It remains to show that \(\theta\notin\ell_\infty(\Gammafin)\) almost surely.
Let \(B_i=\mathbf 1_{\{\xi_i>1\}}\).  By the weighted strong law of large
numbers, or equivalently by summation by parts from the ordinary strong law,
\[
  \frac{\sum_{i=1}^N B_i/i}{\sum_{i=1}^N 1/i}
  \longrightarrow \Prob(\xi_1>1)>0
  \qquad\hbox{almost surely}.
\]
Thus \(\sum_i B_i/i=\infty\) almost surely.

For each \(N\), set \(F_N=\{1\le i\le N:\xi_i>1\}\).  This is a finite subset
of \(\mathbb N\), hence \(F_N\in\Gammafin\), and
\[
  \norm{\theta}_\infty
  \ge \theta(F_N)
  =\sum_{\substack{1\le i\le N\\ \xi_i>1}}\frac{\xi_i}{i}
  \ge \sum_{i=1}^N \frac{B_i}{i}.
\]
The right-hand side is unbounded in \(N\) almost surely.  Therefore
\(\theta\notin U\) almost surely.
\end{proof}

\begin{corollary}[Isomorphic transfer]
Let \(V\) be any Banach space linearly isomorphic to \(\ell_\infty\).  Then
there is a positive nuclear covariance kernel in
\(V\widehat{\otimes}_\pi V\) for which all fixed predual pairings are finite
almost surely, but the associated Gaussian sample is not in \(V\) almost
surely.
\end{corollary}

\begin{proof}
Let \(T:\ell_\infty(\Gammafin)\to V\) be an isomorphism and use the vectors
\(u_i\) from the theorem.  Define
\[
  v_i=\frac{Tu_i}{\norm{Tu_i}_V},\qquad \lambda_i=i^{-2}.
\]
Then \(\norm{v_i}_V=1\) and \(\sum_i\lambda_i<\infty\), so
\(\sum_i \lambda_i v_i\otimes v_i\) is a positive nuclear kernel in
\(V\widehat{\otimes}_\pi V\).  If the corresponding Gaussian series belonged
to \(V\), applying \(T^{-1}\) would give a bounded element of
\(\ell_\infty(\Gammafin)\) with coordinates
\[
  F\mapsto \sum_{i\in F}\frac{\xi_i}{i\,\norm{Tu_i}_V}.
\]
Since \(\norm{Tu_i}_V\le \norm{T}\), the same choice
\(F_N=\{1\le i\le N:\xi_i>1\}\) gives coordinates at least
\(\norm{T}^{-1}\sum_{i\le N}B_i/i\), which are unbounded almost surely.
\end{proof}

\section*{Consequence for the source question}

The source paper itself recalls that, under its Theorem 8.49 hypothesis, the
spaces \(\operatorname{lip}(X^\alpha)\), \(\operatorname{AE}(X^\alpha)\), and
\(\operatorname{Lip}(X^\alpha)\) are linearly homeomorphic to \(c_0\),
\(\ell_1\), and \(\ell_\infty\), respectively.  The corollary therefore gives
a counterexample in those Lipschitz-space instances: a second-predual
structure does not by itself upgrade the source paper's weak-star conclusion
to \(\theta\in U\) almost surely.

\section*{Limitations}

This packet does not rule out stronger positive theorems for special covariance
kernels, for example the exponential kernels studied in the source paper, or
for covariance operators satisfying the appropriate Banach-space Gaussian
radonifying condition.  It only refutes the broad strengthening based solely on
the existence of a second predual.

\section*{Novelty and literature check}

Before promotion, the cheap run indexes were searched for arXiv:2203.04650 and
for the phrases ``Gaussian random fields'', ``non-separable Banach'', ``second
predual'', ``nuclear covariance'', ``\(\ell_\infty\)'', and close variants.
No duplicate packet or attempt was found.

External web searches were run for the exact source title, for
``Gaussian random fields non-separable Banach second predual theta in U'', and
for variants of ``nuclear covariance Gaussian series \(\ell_\infty\) not
bounded''.  The only exact source-title hit was the arXiv paper itself, and no
later paper explicitly answering the Outlook question was found in this
bounded search.  The construction is elementary and should still be reviewed
against the broader Gaussian-measures literature, where the general distinction
between nuclear covariance and radonifying sample operators is classical.

\section*{References}

\begin{itemize}
\item Y. Korolev, J. Latz, C.-B. Schoenlieb, \emph{Gaussian random fields on
non-separable Banach spaces}, arXiv:2203.04650.
\item N. Weaver, \emph{Lipschitz Algebras}, second edition, World Scientific,
2018.  Used here only through Theorem 8.49 as quoted in the source paper.
\end{itemize}

\end{document}
